\subsubsection{Common factors, global phase and relative phase}\label{sec:common-factors-global-phase-and-relative-phase}

Before beginning this section, it is important to review the difference between global and relative phase and the meaning of common factors.

\highspace
Let's consider the following quantum state:
\begin{equation*}
    \ket{\psi} = \sum_j \alpha_j \ket{j}
\end{equation*}
Where $\alpha_j$ are complex numbers and $\ket{j}$ are the basis states. A \textbf{scalar} $c \neq 0$ \textbf{multiplying every amplitude} is a \definition{Common Factor}:
\begin{equation}
    \ket{\phi} = c \ket{\psi} = \sum_j c \alpha_j \ket{j}
\end{equation}
Where $\ket{\phi}$ is the new quantum state. However, we must distinguish three different situations where the common factor $c$ can be applied to the quantum state $\ket{\psi}$:
\begin{itemize}
    \item \definition{Normalization Factor}. Consider the famous quantum state:
    \begin{equation*}
        \ket{+} = \frac{1}{\sqrt{2}}\left(\ket{0} + \ket{1}\right)
    \end{equation*}
    The factor $\frac{1}{\sqrt{2}}$ is necessary because it ensures that the state is normalized:
    \begin{equation*}
        \ket{+} = \underbrace{\frac{1}{\sqrt{2}}}_{\alpha} \ket{0} + \underbrace{\frac{1}{\sqrt{2}}}_{\beta} \ket{1}, \quad \text{where } |\alpha|^2 + |\beta|^2 = 1
    \end{equation*}
    If we remove it, we would have:
    \begin{equation*}
        \ket{0} + \ket{1}, \quad \text{where } |1|^2 + |1|^2 = 2
    \end{equation*}
    The state is not normalized anymore, and it cannot be used in quantum computing. Therefore, a \textbf{normalization factor cannot simply be discarded}.


    \item \definition{Global Phase}. A global phase is a common factor of the form $e^{i \gamma}$, whose modulus is always equal to 1:
    \begin{equation*}
        \left| e^{i \gamma} \right| = 1
    \end{equation*}
    Since the global phase does \textbf{not change the probabilities of measuring the quantum state}, it can be \textbf{discarded}.

    \begin{examplebox}[: Global Phase]
        For example, consider the following quantum state:
        \begin{equation*}
            \ket{\phi} = \frac{i}{\sqrt{2}} \left(\ket{0} + \ket{1}\right)
        \end{equation*}
        Using the Euler's formula, we can rewrite the factor $i$ as:
        \begin{equation*}
            e^{i \theta} = \underbrace{\cos(\theta)}_{\text{must be 0}} + i \underbrace{\sin(\theta)}_{\text{must be 1}} \quad \Rightarrow \quad e^{i \frac{\pi}{2}} = i
        \end{equation*}
        Therefore, we can rewrite the quantum state as:
        \begin{equation*}
            \ket{\phi} = e^{i \frac{\pi}{2}} \frac{1}{\sqrt{2}} \left(\ket{0} + \ket{1}\right) = e^{i \frac{\pi}{2}} \ket{+}
        \end{equation*}
        The two vectors therefore represent the same physical state:
        \begin{equation*}
            \ket{\phi} \sim \ket{+}
        \end{equation*}
        Where $\sim$ means ``\emph{physically equivalent}''. In fact, the global phase can be discarded, and we can write:
        \begin{equation*}
            \ket{\phi} \sim \ket{+} = \frac{1}{\sqrt{2}} \left(\ket{0} + \ket{1}\right)
        \end{equation*}
        Because the measurement probabilities confirm that:
        \begin{equation*}
            P_{\phi}(0) = P_{+}(0) = \left|\frac{i}{\sqrt{2}}\right|^{2} = \frac{1}{2} \quad \text{and} \quad P_{\phi}(1) = P_{+}(1) = \left|\frac{i}{\sqrt{2}}\right|^{2} = \frac{1}{2}
        \end{equation*}
        Thus, the quantum state $\ket{\psi}$ and $e^{i \gamma} \ket{\psi}$ represent the same physical state, and the global phase can be discarded.
    \end{examplebox}


    \item \definition{Relative Phase}. Consider the following quantum states:
    \begin{equation*}
        \ket{+} = \frac{1}{\sqrt{2}} \left(\ket{0} + \ket{1}\right) \quad \text{and} \quad \ket{-} = \frac{1}{\sqrt{2}} \left(\ket{0} - \ket{1}\right)
    \end{equation*}
    The minus sign in the second state multiplieds only the amplitude of the $\ket{1}$ basis state. It is therefore a \definition{Relative Phase}:
    \begin{equation*}
        -1 = e^{i \pi}
    \end{equation*}
    We \textbf{cannot remove} it as a common factor \hl{because it is \textbf{not applied to the entire state}.} Although both states have the same computational-basis probabilities:
    \begin{equation*}
        P_{+}(0) = P_{-}(0) = \left|\frac{1}{\sqrt{2}}\right|^{2} = \frac{1}{2} \quad \text{and} \quad P_{+}(1) = P_{-}(1) = \left|\frac{1}{\sqrt{2}}\right|^{2} = \frac{1}{2}
    \end{equation*}
    The two states are \textbf{not physically equivalent}:
    \begin{equation*}
        \ket{+} \nsim \ket{-}
    \end{equation*}
    Because the relative phase affects the interference of the two states, for example, when we apply the Hadamard gate to both states:
    \begin{equation*}
        H\ket{+} = \ket{0} \quad \text{but} \quad H\ket{-} = \ket{1}
    \end{equation*}
    Thus, relative phases are \textbf{physically relevant}.
\end{itemize}
In general, suppose we have two quantum states:
\begin{equation*}
    \ket{\psi} = \sum_j \alpha_j \ket{j}, \qquad \ket{\phi} = \sum_j \beta_j \ket{j}
\end{equation*}
If both states are normalized, they represent the same physical state exactly when there is a single number $c$ such that:
\begin{equation*}
    \beta_j = c \alpha_j \quad \forall j, \, \left|c\right| = 1
\end{equation*}
In other words, all corresponding nonzero amplitudes must have the same ratio:
\begin{equation*}
    \frac{\beta_0}{\alpha_0} = \frac{\beta_1}{\alpha_1} = \cdots = e^{i \gamma}
\end{equation*}
If the ratio is not the same for all amplitudes, then the two states are \textbf{not physically equivalent} because they have different relative phases.

\highspace
\exercise Determine whether the following states represent the same physical state:
\begin{equation*}
    \ket{\psi} = \frac{1}{\sqrt{2}} \left(\ket{0} + \ket{1}\right), \qquad \ket{\phi} = \frac{1+i}{2} \left(\ket{0} + \ket{1}\right)
\end{equation*}
\solution We need to determine whether there is a single number $c$ such that:
\begin{equation*}
    \frac{1+i}{2} = c \cdot \frac{1}{\sqrt{2}}
\end{equation*}
Or in other words, we need to determine whether $\frac{1+i}{2}$ differs from $\frac{1}{\sqrt{2}}$ by a global phase.
\begin{enumerate}
    \item First of all, check that \textbf{both states are normalized}:
    \begin{equation*}
        \left|\frac{1}{\sqrt{2}}\right|^2 + \left|\frac{1}{\sqrt{2}}\right|^2 = \frac{1}{2} + \frac{1}{2} = 1,
    \end{equation*}
    \begin{equation*}
        \begin{aligned}
            \left|\frac{1+i}{2}\right|^2 &= \frac{\left|1+i\right|^{2}}{\left|2\right|^{2}} \\
            &= \frac{\left|1+i\right|^{2}}{4} \\
            &\downarrow \text{ Remark: } \left|z\right| = \sqrt{a^2 + b^2} \text{ and } \left|z\right|^{2} = a^2 + b^2 \text{ for } z = a + ib \\
            &\downarrow \text{ For } 1+i, \text{ we have } a = 1, b = 1 \\
            &= \frac{1^2 + 1^2}{4} \\
            &= \frac{2}{4} \\
            &= \frac{1}{2} \Rightarrow \left|\frac{1+i}{2}\right|^2 + \left|\frac{1+i}{2}\right|^2 = \frac{1}{2} + \frac{1}{2} = 1
        \end{aligned}
    \end{equation*}

    
    \item Now, we can \textbf{find the factor connecting the two states}. In general, we want to know if one state can be obtained by multiplying the entirety of the other state by one number $c$:
    \begin{equation*}
        \phi = c \psi
    \end{equation*}
    Let's substitute the two states:
    \begin{equation*}
        \frac{1+i}{2} \left(\ket{0} + \ket{1}\right) = c \cdot \frac{1}{\sqrt{2}} \left(\ket{0} + \ket{1}\right)
    \end{equation*}
    We can divide both sides by $\left(\ket{0} + \ket{1}\right)$ to remove the basis states:
    \begin{equation*}
        \frac{1+i}{2} = c \cdot \frac{1}{\sqrt{2}}
    \end{equation*}


    \item Now, we can solve for $c$:
    \begin{equation*}
        \begin{aligned}
            \frac{1+i}{2} &= c \cdot \frac{1}{\sqrt{2}} \\
            \sqrt{2} \cdot \frac{1+i}{2} &= c \\
            c &= \sqrt{2} \cdot \frac{1+i}{2}
        \end{aligned}
    \end{equation*}
    And \textbf{verify} that multiplying $\psi$ by $c$ gives $\phi$ (we need to check that $c$ is a global phase):
    \begin{align*}
        \sqrt{2} \cdot \frac{1+i}{2} \ket{\psi} &= \sqrt{2} \cdot \frac{1+i}{2} \cdot \left[\frac{1}{\sqrt{2}}\left(\ket{0} + \ket{1}\right)\right] \\
        &= \sqrt{2} \cdot \frac{1+i}{2} \cdot \frac{1}{\sqrt{2}} \left(\ket{0} + \ket{1}\right) \\
        &= \cancel{\sqrt{2}} \cdot \frac{1+i}{2} \cdot \frac{1}{\cancel{\sqrt{2}}} \left(\ket{0} + \ket{1}\right) \\
        &= \frac{1+i}{2} \left(\ket{0} + \ket{1}\right) \\
        &= \ket{\phi}
    \end{align*}
    So we have found one common factor multiplying the entire state, but this does not mean that the two states are physically equivalent! We need to check that the factor is a global phase, i.e. that its modulus is equal to 1:
    \begin{align*}
        \left|c\right| &= \left|\sqrt{2} \cdot \frac{1+i}{2}\right| \\
        &= \left|\sqrt{2}\right| \cdot \left|\frac{1+i}{2}\right| \\
        &= \sqrt{2} \cdot \frac{\left|1+i\right|}{\left|2\right|} \\
        &= \sqrt{2} \cdot \frac{\left|1+i\right|}{2} \\
        &\downarrow \text{ Remark: } \left|z\right| = \sqrt{a^2 + b^2} \text{ and } \left|z\right|^{2} = a^2 + b^2 \text{ for } z = a + ib \\
        &\downarrow \text{ For } 1+i, \text{ we have } a = 1, b = 1 \\
        &= \sqrt{2} \cdot \frac{\sqrt{1^2 + 1^2}}{2} \\
        &= \sqrt{2} \cdot \frac{\sqrt{2}}{2} \\
        &= \frac{2}{2} = 1
    \end{align*}
\end{enumerate}
Thus, $c$ \textbf{changes only the phase}; it does \textbf{not change the magnitude of the amplitudes}. Therefore, $c$ is a global phase:
\begin{equation*}
    \ket{\phi} = c \ket{\psi}, \quad \text{where } \left|c\right| = 1
\end{equation*}
And the states are physically equivalent: $\ket{\phi} \sim \ket{\psi}$.

\highspace
\exercise Determine whether these states represent the same physical state:
\begin{equation*}
    \ket{\psi} = \frac{1}{2} \left(i\ket{00} + \ket{01} - i\ket{10} + \ket{11}\right), \quad
    \ket{\phi} = \frac{1}{2} \left(\ket{00} - i\ket{01} - \ket{10} - i\ket{11}\right)
\end{equation*}
\solution We need to determine whether there is a single number $c$ such that:
\begin{equation*}
    \begin{aligned}
        \ket{\psi} &= c \ket{\phi} \\
        \frac{1}{2} \left(i\ket{00} + \ket{01} - i\ket{10} + \ket{11}\right) &= c \cdot \frac{1}{2} \left(\ket{00} - i\ket{01} - \ket{10} - i\ket{11}\right)
    \end{aligned}
\end{equation*}
Let's try to find the factor $c$ by comparing the amplitudes of the basis states:
\begin{enumerate}
    \item First of all, check that \textbf{both states are normalized}. We start with $\ket{\psi}$:
    \begin{align*}
        \underbrace{\left|\frac{1}{2} \cdot i\right|^{2}}_{\ket{00}} + \underbrace{\left|\frac{1}{2}\right|^{2}}_{\ket{01}} + \underbrace{\left|\frac{1}{2} \cdot (-i)\right|^{2}}_{\ket{10}} + \underbrace{\left|\frac{1}{2}\right|^{2}}_{\ket{11}} &= 1 \\
        \left(\sqrt{0^2 + \left(\frac{1}{2}\right)^2}\right)^{2} + \frac{1}{4} + \left(\sqrt{0^2 + \left(-\frac{1}{2}\right)^2}\right)^{2} + \frac{1}{4} &= 1 \\
        \frac{1}{4} + \frac{1}{4} + \frac{1}{4} + \frac{1}{4} &= 1 \\
        1 &= 1
    \end{align*}
    And now for $\ket{\phi}$:
    \begin{align*}
        \underbrace{\left|\frac{1}{2}\right|^{2}}_{\ket{00}} + \underbrace{\left|\frac{1}{2} \cdot (-i)\right|^{2}}_{\ket{01}} + \underbrace{\left|\frac{1}{2} \cdot (-1)\right|^{2}}_{\ket{10}} + \underbrace{\left|\frac{1}{2} \cdot (-i)\right|^{2}}_{\ket{11}} &= 1 \\
        \frac{1}{4} + \left(\sqrt{0^2 + \left(-\frac{1}{2}\right)^2}\right)^{2} + \frac{1}{4} + \left(\sqrt{0^2 + \left(-\frac{1}{2}\right)^2}\right)^{2} &= 1 \\
        \frac{1}{4} + \frac{1}{4} + \frac{1}{4} + \frac{1}{4} &= 1 \\
        1 &= 1
    \end{align*}


    \item Now, we can \textbf{find the factor connecting the two states}, but we need to check that the factor is the same for all amplitudes.
    \begin{itemize}
        \item For the $\ket{00}$ basis state, we have:
        \begin{equation*}
            \frac{1}{2} \cdot i = c \cdot \frac{1}{2} \cdot 1 \quad \Rightarrow \quad c = i
        \end{equation*}
        \item For the $\ket{01}$ basis state, we have:
        \begin{equation*}
                \frac{1}{2} = c \cdot \frac{1}{2} \cdot (-i) \quad \Rightarrow \quad 1 = c \cdot (-i) \quad \Rightarrow \quad c = -\frac{1}{i}
        \end{equation*}
        \item For the $\ket{10}$ basis state, we have:
        \begin{equation*}
            -\frac{1}{2} \cdot i = c \cdot \frac{1}{2} \cdot (-1) \quad \Rightarrow \quad c = i
        \end{equation*}
        \item For the $\ket{11}$ basis state, we have:
        \begin{equation*}
            \frac{1}{2} = c \cdot \frac{1}{2} \cdot (-i) \quad \Rightarrow \quad c = -\frac{1}{i}
        \end{equation*}
    \end{itemize}
    We can simplify each $c = -\frac{1}{i}$:
    \begin{equation*}
        -\frac{1}{i} = -\frac{1}{i} \cdot \underbrace{\frac{i}{i}}_{1} = -\frac{i}{i^2} = -\frac{i}{-1} = i
    \end{equation*}
    Thus, we have found that the factor $c$ is the same for all amplitudes:
    \begin{equation*}
        c = i
    \end{equation*}
    

    \item The previous step shows that the two states are related by a common factor $c = i$. However, we need to check that $c$ is a global phase, i.e. that its modulus is equal to 1:
    \begin{equation*}
        \left|c\right| = \left|i\right| = \sqrt{0^2 + 1^2} = \sqrt{1} = 1
    \end{equation*}
\end{enumerate}
A global phase means that there exists \textbf{\underline{one} single complex number} $c$, with modulus equal to 1, such that:
\begin{equation*}
    \ket{\psi} = c \ket{\phi}
\end{equation*}